\section*{摘要}
本实验分析完整编译工具链，并对两个互补的 SysY 示例分别手写 LLVM IR 与 RISC-V 汇编。工具链部分比较 GCC/Clang 的 O0、O2 与 O2-g 六种构建，60 次运行均符合预期；Clang 的向量化伴随代码增大，而两种编译器各自加入调试信息后代码段内容一致。语义等价部分覆盖短路副作用、循环跳转、二维数组、函数、作用域及单精度浮点，链接公开的 SysY 运行库，在 QEMU 中完成 160 个用例、三种实现共 480 次执行验证，并结合 CFG、SSA 前驱与函数调用分析其等价机制。进阶部分使用官方 AscendNPU IR 1.1.0，在 x86\_64 主机上对面向 Ascend 设备的 VecAdd IR 完成六次处理，分析 AIV 核推导、UB 地址复用、跨流水同步和模板调用转换。由于本地缺少 hivmc/CANN 与昇腾设备，设备目标文件生成和设备端执行尚未完成。

\noindent\textbf{关键词：} 编译工具链；SysY；LLVM IR；RISC-V；SSA；AscendNPU IR
